% ============================================================
% FULL METHOD SECTION — Gina (pipeline) + Chiebuka (improvements)
% ============================================================


% ────────────────────────────────────────────────────────────
\section{Method}
% ────────────────────────────────────────────────────────────

\subsection{System Pipeline}

The system is structured as a sequential four-stage pipeline,
taking raw Yelp data as input and producing a ranked list of
businesses as output.

\paragraph{Stage 1 — Data Loading.}
The Yelp Academic Dataset is stored in JSON Lines format, with
one business record per line.
We load up to 5,000 records and retain four fields:
business name, star rating (1–5), review count, and categories
(a comma-separated string of tags such as ``restaurants, chinese, szechuan'').
All other fields (address, hours, attributes, etc.) are discarded
to keep the working set compact.

\paragraph{Stage 2 — Preprocessing.}
Raw records often contain missing or malformed values.
We apply three cleaning steps:
(1) drop any record missing one or more of the four retained fields;
(2) coerce star ratings and review counts to numeric types,
    dropping rows that cannot be converted;
(3) lowercase and strip whitespace from the categories string so
    that keyword comparisons are case-insensitive.
After preprocessing, the working dataset contains 3,943 businesses.

\paragraph{Stage 3 — Filtering.}
Before ranking, a hard filter removes businesses below a
user-specified minimum star rating (default: 3.0).
This acts as a baseline quality gate, ensuring that no result
with an unacceptably low rating is ever returned regardless of
its keyword relevance or review count.
Filtering is intentionally kept simple so that the ranking
function, not the filter, is responsible for ordering results.

\paragraph{Stage 4 — Ranking and Output.}
The filtered set is passed to a ranking function (described in
Sections~\ref{sec:baseline} and~\ref{sec:improved}) which assigns
a scalar score to each business and returns the top-$K$.
Results are written to a CSV file for reproducibility and
rendered as a horizontal bar chart for visual inspection.
The full pipeline is invoked through a single entry point:
\texttt{python -m src.query}.

\subsection{Baseline Ranking Model}
\label{sec:baseline}

The baseline system ranks businesses using a weighted linear combination
of three normalized signals:

\begin{equation}
  s_{\text{base}}(b) \;=\;
    w_r \cdot \frac{\text{stars}(b)}{5}
    \;+\; w_v \cdot \min\!\left(\frac{\text{reviews}(b)}{1000},\;1\right)
    \;+\; w_k \cdot \text{KM}(b,\,Q)
  \label{eq:baseline}
\end{equation}

\noindent where the default weights are $w_r = 0.5$, $w_v = 0.3$,
$w_k = 0.2$, and the keyword match score is

\begin{equation}
  \text{KM}(b, Q) \;=\;
    \frac{1}{|Q|}\sum_{q \in Q} \mathbf{1}\!\left[q \in \text{cat}(b)\right]
  \label{eq:km_baseline}
\end{equation}

\noindent with $\text{cat}(b)$ denoting the lowercased category string of
business $b$ and $Q$ the set of user-provided query keywords.
All scores are bounded in $[0, 1]$ and weights sum to 1.

\subsection{Improved Ranking Model}
\label{sec:improved}

We identify three weaknesses in the baseline and address each:

\paragraph{1. Review count normalization.}
The baseline clips review count at 1,000, treating a business with
1,001 reviews identically to one with 100,000.
Popularity on Yelp follows a heavy-tailed distribution, so a
logarithmic normalization is more appropriate:

\begin{equation}
  \phi(b) \;=\;
    \frac{\ln\!\left(1 + \text{reviews}(b)\right)}
         {\ln\!\left(1 + \max_{b' \in \mathcal{B}}\,\text{reviews}(b')\right)}
  \label{eq:log_norm}
\end{equation}

\noindent This keeps scores in $[0,1]$, compresses the long tail, and
preserves relative ordering across the full review-count range.

\paragraph{2. Keyword matching scope.}
The baseline matches query keywords only against the \texttt{categories}
field.  Many businesses encode relevant terms in their name (e.g.\
\textit{Han Dynasty}) rather than only in categories.  We extend
matching to include the business name with a discounted credit:

\begin{equation}
  \text{KM}'(b, Q) \;=\;
    \frac{1}{|Q|}\sum_{q \in Q}
    \begin{cases}
      1.0 & \text{if } q \in \text{cat}(b) \\
      0.5 & \text{if } q \in \text{name}(b) \text{ and } q \notin \text{cat}(b) \\
      0   & \text{otherwise}
    \end{cases}
  \label{eq:km_improved}
\end{equation}

\paragraph{3. Keyword weight.}
Because query relevance should be a primary ranking signal, we raise
$w_k$ from 0.2 to 0.3 and redistribute the remaining weight
proportionally to preserve the original rating-to-review ratio:

\begin{equation}
  w_r = 0.4,\quad w_v = 0.3,\quad w_k = 0.3,\quad w_r + w_v + w_k = 1
  \label{eq:weights}
\end{equation}

\paragraph{Improved scoring function.}
Substituting the above into \eqref{eq:baseline}:

\begin{equation}
  s_{\text{imp}}(b) \;=\;
    0.4\cdot\frac{\text{stars}(b)}{5}
    \;+\; 0.3\cdot\phi(b)
    \;+\; 0.3\cdot\text{KM}'(b,\,Q)
  \label{eq:improved}
\end{equation}

\paragraph{Tie-breaking.}
When two businesses share the same score, we break ties by
$\text{reviews}(b)$ descending, giving preference to the more
widely reviewed option as a secondary popularity signal.

\paragraph{Top-$K$ selection.}
Given dataset $\mathcal{B}$ and query $Q$, the output is

\begin{equation}
  \mathcal{R}^* \;=\;
    \underset{\mathcal{R} \subseteq \mathcal{B},\,|\mathcal{R}|=K}{\arg\max}
    \sum_{b \in \mathcal{R}} s_{\text{imp}}(b)
\end{equation}

\noindent which is solved by sorting all businesses by score and
returning the top $K$ — an $O(n \log n)$ operation for dataset size $n$.


% ────────────────────────────────────────────────────────────
\section{Experimental Results}
% ────────────────────────────────────────────────────────────

\subsection{Experimental Setup}

All experiments use the Yelp Academic Dataset (business subset),
loaded as JSON-lines and preprocessed to retain businesses with
non-null name, category, star rating, and review count.
After filtering to a minimum rating of 3.0 stars, the working set
contains 3,943 businesses.
We fix $K = 5$ unless stated otherwise, and evaluate three queries:
\textit{chinese + spicy}, \textit{pizza + italian},
and \textit{coffee + brunch}.

\subsection{Evaluation Metric: Precision@$K$}

We define relevance as binary: business $b$ is relevant to query $Q$ if

\begin{equation}
  \text{rel}(b, Q) \;=\; \mathbf{1}\!\left[\exists\, q \in Q : q \in \text{cat}(b)\right]
\end{equation}

\noindent Precision at rank $K$ is then

\begin{equation}
  P@K \;=\; \frac{1}{K}\sum_{i=1}^{K} \text{rel}(b_i,\, Q)
  \label{eq:precision}
\end{equation}

\noindent where $b_1, \ldots, b_K$ are the ranked results.
This metric directly measures whether the returned businesses
match the user's stated query intent.

\subsection{Experiment 1: Multi-Query Comparison}

Table~\ref{tab:multiquery} shows $P@5$ for the baseline and improved
ranker across all three queries (Figure~\ref{fig:exp1}).

\begin{table}[h]
\centering
\begin{tabular}{lcc}
\hline
\textbf{Query} & \textbf{Baseline $P@5$} & \textbf{Improved $P@5$} \\
\hline
chinese + spicy  & 0.20 & \textbf{1.00} \\
pizza + italian  & 0.40 & \textbf{1.00} \\
coffee + brunch  & 1.00 & \textbf{1.00} \\
\hline
\end{tabular}
\caption{Precision@5 comparison across three queries.}
\label{tab:multiquery}
\end{table}

The baseline fails most severely on \textit{chinese + spicy}: its top-3
results are non-Chinese restaurants (\textit{Prep \& Pastry},
\textit{District Donuts}, \textit{Katie's Restaurant \& Bar}) that
achieve high scores purely through star ratings and review counts.
The improved model ranks \textit{Han Dynasty} first on this query,
demonstrating that upweighting keyword relevance corrects for
popularity bias.
The \textit{coffee + brunch} query achieves $P@5 = 1.0$ under both
methods because coffee and brunch are common categories, providing
sufficient keyword signal even at low weight.

\subsection{Experiment 2: Precision@$K$ Across Ranks}

Table~\ref{tab:precatk} reports $P@K$ for $K \in \{1, 3, 5, 10\}$
(Figure~\ref{fig:exp2}).

\begin{table}[h]
\centering
\begin{tabular}{llcccc}
\hline
\textbf{Query} & \textbf{Method} & $P@1$ & $P@3$ & $P@5$ & $P@10$ \\
\hline
chinese + spicy & Baseline & 0.00 & 0.00 & 0.20 & 0.10 \\
                & Improved & \textbf{1.00} & \textbf{1.00} & \textbf{1.00} & \textbf{1.00} \\
\hline
pizza + italian & Baseline & 1.00 & 0.67 & 0.40 & 0.40 \\
                & Improved & \textbf{1.00} & \textbf{1.00} & \textbf{1.00} & \textbf{1.00} \\
\hline
coffee + brunch & Baseline & 1.00 & 1.00 & 1.00 & 0.90 \\
                & Improved & \textbf{1.00} & \textbf{1.00} & \textbf{1.00} & \textbf{1.00} \\
\hline
\end{tabular}
\caption{Precision@$K$ for both rankers across queries and rank cutoffs.}
\label{tab:precatk}
\end{table}

The improved ranker achieves perfect precision at every cutoff and
every query.
Baseline precision degrades as $K$ grows because high-review-count
irrelevant businesses accumulate in lower ranks.
This confirms that the popularity bias in the baseline is a
structural limitation, not a marginal effect.

\subsection{Experiment 3: Keyword Weight Sensitivity}

To justify the choice $w_k = 0.3$, we sweep $w_k$ from 0 to 0.6,
redistributing the remaining weight between $w_r$ and $w_v$ in a
fixed 5:3 ratio (Figure~\ref{fig:exp3}).

Key observations:
\begin{itemize}
  \item At $w_k = 0$ (pure popularity), $P@5 = 0.00$ for
        \textit{chinese + spicy} — the model returns no relevant results.
  \item Precision rises to 0.40 at $w_k = 0.20$ (baseline weight)
        and reaches 1.00 at $w_k \geq 0.25$.
  \item The mean score of irrelevant results in the top-$K$ drops to
        zero at $w_k = 0.25$, meaning no irrelevant business can
        outrank a keyword-matching one past this threshold.
  \item $w_k = 0.30$ sits comfortably past the transition point,
        providing a margin against edge cases while keeping
        rating and review signals influential.
\end{itemize}

\subsection{Discussion \& Future Work}

The results demonstrate that a modest increase in keyword weight
eliminates the popularity bias that plagues the baseline while
preserving the value of quality and engagement signals.
The log normalization of review counts prevents a business with
10,000 reviews from dominating one with 1,000 reviews, allowing
truly relevant but moderately popular businesses to surface.

One limitation is that our relevance metric (keyword in category)
is imperfect: \textit{Dawna Ara, DACM, LAc} (an acupuncture clinic)
ranks second on \textit{chinese + spicy} because ``traditional chinese
medicine'' appears in its categories.
Future work could address this by (1) restricting matches to a
business-type filter (e.g.\ restaurants only), (2) incorporating
semantic similarity via word embeddings, or (3) learning weights
from user click data rather than setting them manually.
